\section{The original point moment and the growing-cone phase}
\label{S:original-point-transfer}

Retain all original covariances and amplitudes. Write $a = \alpha_0$, $\kappa =
(4-\alpha-a)/(3-\alpha) > 1$. Fix the auxiliary parameters of the preceding
Gaussian and product sections once, in their strict admissible gaps. For $0 <
\delta < 1$ put
\[
  L=\delta^{-1},\quad R=L+2,\quad T_\delta=\delta^{-1/\kappa},
  \quad g=\delta^{(1-a)/2},\quad \Gamma_\delta=g^2\Gamma .
\]
Use the same real even root $\eta$, set $\psi = \eta(\cdot/R)$, $m = \psi^2$, $z
= (\psi,0) \in E$, $D = B_1(0)$ and $v_d = |D| > 0$. Let $W_{\rm fast} = W\circ
T_{\delta,0}$ on the given slow probability space, and let $P_\delta$ be the
physical $E$ endpoint at time $L$ driven by $gM_R(W_{\rm fast})$, with physical
noise mask $m$ once. The datum $\psi$ is not the auxiliary cutoff used to extend
the operator path to constant columns on $F$.

\begin{lemma}[Restriction and fixed-parameter finiteness]
  \label{L:point-actual-finiteness}
  For nonzero covariance, the $E$ flow and its source flows are the closed
  isometric restrictions of their $F$ extensions. Their operator norms and
  differences are measurable, and
  \begin{equation}\label{E:point-fixed-finiteness}
    1\le\E\|P_\delta\|_{\mathcal{L}(E)}^2<\infty .
  \end{equation}
\end{lemma}

\begin{proof}
  We first identify the $E$ flow as a closed isometric restriction of the
  extended phase flow. We then apply the fixed-parameter Gaussian tube
  estimate and compact endpoint lower bound to this exact original endpoint.


  \medskip
  \noindent\textit{Step 1.~}We identify restrictions and measurable endpoint
  norms.
  Let $A_F = gA_{F,R}(W_{\rm fast})$ be the scaled full operator path. Linearity
  and Proposition~\ref{P:constant-phase-driver} give $A_Fi = iD_E$, with $D_E =
  gM_R(W_{\rm fast})$, and $S_F(t)i = iS_E(t)$. The extension uses a cutoff
  $\widetilde{\psi} = 1$ near its larger localization, not the datum $\psi$.
  Applying $i$ to the state Young sums and using uniqueness gives $U_F(t,s)i =
  iU_E(t,s)$, then $P_F(t,s)i = iP_E(t,s)$. The same argument applies to
  sources. Closedness and isometry imply $\|P_E\| \le \|P_F\|$ and the analogous
  difference bound; density of $E$ in $F$ is neither true nor used. The
  continuous Young map on its separable path span, followed by bounded physical
  conjugation and restriction, gives measurable endpoint norms and uniform
  time-triangle difference norms. No operator-norm continuity of the free group
  is assumed. Young Neumann pieces and the free inverse also give invertibility.


  \medskip
  \noindent\textit{Step 2.~}We prove finiteness and the lower endpoint norm
  bound.
  We check a direct finiteness input, without a fixed-mask asymptotic rate. Take
  $B = 1$ in Propositions~\ref{P:tube-actual-law}
  and~\ref{P:tube-coarse-envelope} and Lemma~\ref{L:tube-rayleigh-tail}. Every
  $0 < \delta < 1$ satisfies their clock condition. The global endpoint
  in~\eqref{E:tube-global-endpoint} is exactly this $P_\delta$: the same given
  $W_{\rm fast}$, amplitude $g$, root $\eta$, one mask $m$, massless $E$ phase,
  and horizon $\delta^{-1}$. We use the same admissible $\theta$ and operator
  path version. Explicit row identities, rational times and Young uniqueness
  identify versions if a smaller common admissible exponent is used. Thus, the
  finite radial bound applies to this endpoint, not merely an equal-law copy. At
  these fixed parameters the operator path $gM_R(W_{\rm fast})$ is a
  compact-valued $C^\theta([0,L];\mathcal{K}(E))$ path starting at zero on its
  path event. Lemma~\ref{L:symplectic-free} gives $\|P_\delta\| \ge \|S_E(L)\|
  \ge 1$ almost surely, proving~\eqref{E:point-fixed-finiteness}. This use
  requires no $F$-norm rate along growing radii. Source convergence and source
  uniform integrability are checked next, separately from this original endpoint
  finiteness. This proves Lemma~\ref{L:point-actual-finiteness}.
\end{proof}

\begin{lemma}[The same regularized fast noises preserving the given noise]\label{L:point-fast-sources}
  At each fixed $\delta$, take one regularized slow noise preserving the given
  noise $Y_{\epsilon,\tau}$ and define
  \[
    Y_{\rm fast}(t,x)=\delta^{a/2}Y_{\epsilon,\tau}(\delta t,x).
  \]
  Its physical $E$ flows $P_E^{\epsilon,\tau}$ with coefficient $m gY_{\rm
  fast}$ satisfy
  \begin{equation}\label{E:point-source-flow-limit}
    \E\sup_{0\le s\le t\le L}
    \|P_E^{\epsilon,\tau}(t,s)-P_E(t,s)\|^2\longrightarrow0
    \quad\hbox{as }(\epsilon,\tau)\to(0,0).
  \end{equation}
  The unmasked field $gY_{\rm fast}$ has the canonical source law for covariance
  $\Gamma_\delta$, spatial width $\epsilon$ and temporal width $\tau/\delta$.
\end{lemma}

\begin{proof}
  We pull one regularized slow noise into fast coordinates and transfer its
  operator paths on the original coupling. We identify the ordinary response
  and compute the exact temporal width and covariance amplitude.


  \medskip
  \noindent\textit{Step 1.~}We prove convergence of the source flow at fixed
  dilation.
  Testwise this field is $Y_{\epsilon,\tau}(T_{\delta,0}f)$. The original
  isometry and source domination give source and residual variances bounded by
  $\|f\|_{H_{\rm fast}}^2$, with residual variance tending to zero at fixed
  $\delta$ on every smooth compact test. The Gaussian transfer
  Proposition~\ref{P:source-gaussian-transfer} therefore gives operator-path
  convergence in every finite $L^pC^\theta$, retaining $W_{\rm fast}$ literally.
  Multiplication by fixed $g$ and the bounded $F$ extension preserve it. The
  transfer bound and Proposition~\ref{P:actual-operator-driver} also give
  uniform source $\sqrt p$ path moments at these fixed $\delta,R,L$. The
  moment-series argument yields square-exponential path norms. Young local
  dependence gives probability convergence on ball of operator pathss; the
  growth exponent $1/\theta < 2$ and those tails give uniform flow moments of
  every finite order and uniform integrability. The proof of
  Proposition~\ref{P:constant-gaussian-flow} applies to this very pullback
  family. Restriction from $F$ then proves~\eqref{E:point-source-flow-limit}. No
  uniform-in-$\delta$ source bound or exchange of limits is asserted.


  \medskip
  \noindent\textit{Step 2.~}We identify the source primitive on the same
  states.
  This source is also the ordinary response for all $E$ data. On Schwartz states
  its first primitive has row $gY_{\rm fast}(h_R)$. The source indicator
  approximation and bounded point variances at positive widths identify it with
  the statewise integral $\int_0^t S_E(-r)B(mgY_{\rm fast}(r))S_E(r)\ud r$.
  Fubini uses $\E|ZY| \le \|Z\|_2\|Y\|_2$ when Gaussian testing is needed.
  Countable dense states, rational times and continuity give the full
  ordinary-primitive identity, as in
  Proposition~\ref{P:source-ordinary-response}.
  Lemma~\ref{L:constant-source-solution} and
  Proposition~\ref{P:symplectic-source} then identify the Young equation with
  the ordinary interaction equation on $E$.


  \medskip
  \noindent\textit{Step 3.~}We compute the transformed source law.
  In the temporal kernel $K_\tau(v) = \Gamma_{\rm E}(a)^{-1}
  \int_0^{\tau^{-1}}e^{-\lambda|v|}\lambda^{a-1}\ud\lambda$, substitution
  $\lambda' = \delta\lambda$ gives
  \[
    \delta^aK_\tau(\delta v)=K_{\tau/\delta}(v).
  \]
  The change of variable contributes $\delta^{-1}$ from $d\lambda$ and
  $\delta^{1-a}$ from $\lambda^{a-1}$. These cancel the prefactor $\delta^a$,
  while the upper cutoff becomes $\delta/\tau$. Thus, the new temporal width
  is $\tau/\delta$, with the original temporal covariance amplitude unchanged
  before multiplication by $g^2$.
  Hence, $gY_{\rm fast}$ has covariance
  $K_{\tau/\delta}(t-s)(p_{2\epsilon}*\Gamma_\delta)(x-y)$. The shared-width
  maximum and spatial heat-sum conventions persist for two sources. Positive
  multiplication by $g^2$ preserves every original covariance class, including
  spatial white noise. At fixed $\delta$, $(\epsilon,\tau/\delta) \to (0,0)$.
  The full source point-moment limit in
  Proposition~\ref{P:masked-canonical-transfer} thus applies to this
  single-source law and its continuous version; no newly reconstructed auxiliary
  family for $\Gamma_\delta$ is required. This proves
  Lemma~\ref{L:point-fast-sources}.
\end{proof}

\begin{lemma}[Compact-data source identification]\label{L:point-compact-source}
  For fixed $\delta,\epsilon,\tau$, the displacement $v_{\epsilon,\tau}(t,x) =
  (\operatorname{pr}_1P_E^{\epsilon,\tau}(t,0)z)(x)$ has a jointly measurable
  continuous version and satisfies the ordinary scalar mild equation with
  coefficient $b = mgY_{\rm fast}$ and initial data $(\psi,0)$.
\end{lemma}

\begin{proof}
  At positive source widths we solve in a sufficiently regular Sobolev phase,
  identify that solution with the $E$ flow by uniqueness, and then use Fourier
  embedding and Green testing to recover the continuous scalar mild solution.


  \medskip
  \noindent\textit{Step 1.~}We construct the high-regularity phase solution.
  On the fixed source regularity event, $b$ is smooth in space, has common
  compact support in $B_{2R}$, and is continuous in time in every bounded
  spatial derivative norm on $[0,L]$. Choose integer $k > d/2+1$ and $E_k =
  H^k\times H^{k-1}$. Its Bessel-unitary free matrix is
  \[
    \begin{pmatrix}
      \cos(t|\xi|)&\langle\xi\rangle\sin(t|\xi|)/|\xi|\\
      -|\xi|\sin(t|\xi|)/\langle\xi\rangle&\cos(t|\xi|)
    \end{pmatrix}.
  \]
  The zero sine quotient is $t$; finite-horizon bounds and dominated Fourier
  convergence give a strongly continuous group on $E_k$. Leibniz and the
  derivative-norm equivalence, obtained by expanding $(1+|\xi|^2)^{k-1}$ and
  Parseval, give
  \[
    \|b(t)f\|_{H^{k-1}}
    \le C_k\max_{|j|\le k-1}\|\partial^jb(t)\|_\infty
    \|f\|_{H^{k-1}} .
  \]
  Thus, $B(b(t))$ is a norm-continuous bounded force on $E_k$, with a finite
  uniform sample bound $B_L$. For $Z_0 = S(t)z$ and $Z_{n+1}(t) =
  \int_0^tS(t-r)B(b(r))Z_n(r)\ud r$,
  \[
    \sup_{t\le L}\|Z_n(t)\|_{E_k}
    \le C_L^{n+1}B_L^nL^n\|z\|_{E_k}/n! .
  \]
  This series converges in $C_tE_k$. Iterating a homogeneous Volterra difference
  proves uniqueness. Its continuous embedding into $E$ and ordinary $E$
  uniqueness identify the sum with $P_E^{\epsilon,\tau}(t,0)z$.


  \medskip
  \noindent\textit{Step 2.~}We obtain continuous measurable displacement and
  velocity.
  Fourier Cauchy--Schwarz gives $\|\widehat{f}\|_1 \le \|f\|_{H^k}
  (\int\langle\xi\rangle^{-2k}\ud\xi)^{1/2} < \infty$. Fourier inversion and
  dominated convergence imply continuous bounded displacement, and the same
  holds for velocity since $k-1 > d/2$. The source derivatives on the compact
  support form a measurable $C_tC_x^{k-1}$ random element, by countable
  evaluations and continuity. Picard iteration and local norm estimates make its
  phase sum measurable in $C_tE_k$; the Fourier embedding gives joint
  measurability. These high norms are used only at positive source widths.


  \medskip
  \noindent\textit{Step 3.~}We identify the scalar Green equation.
  The choice $k-1>d/2$ makes both Fourier embeddings available: displacement
  uses $H^k$ and velocity uses $H^{k-1}$. Continuity in the phase norm
  therefore gives continuity of both components as functions of space and time
  at these fixed positive source widths. This higher regularity is not claimed
  for the original rough-noise solution.
  The first component of Duhamel is, distributionally,
  \[
    v(t)=\cos(t\sqrt{-\Delta})\psi+
    \int_0^t G_d(t-r)*[b(r)v(r)]\ud r .
  \]
  Here $G_d(h)$ is the positive Green measure of mass $h$ and support in $B_h$,
  including the sphere measure for $d = 3$. The product $bv$ is continuous
  compactly supported, with finite uniform $L^1$ and supremum bounds. Against a
  compact smooth $\phi$, absolute Fubini is bounded by
  $\int_0^t\|\phi\|_\infty\|b(r)v(r)\|_1(t-r)\ud r < \infty$. Writing $G_d(h)$
  as $h$ times the dilation of its unit-time probability measure proves
  continuity of the convolution term. The free term is continuous by the Fourier
  embedding. Thus, the distributional identity agrees everywhere with the
  pointwise continuous mild versions. No Green density in $d = 3$ or original
  Sobolev point trace is used. This proves Lemma~\ref{L:point-compact-source}.
\end{proof}

The displacement $v_{\epsilon,\tau}$ of Lemma~\ref{L:point-compact-source} has
the masked coefficient $mgY_{\rm fast}$ and the data $(\psi,0)$. The unmasked
source solution has the coefficient $gY_{\rm fast}$ and the data $(1,0)$. We now
show that the two agree on a cone that contains $D$ at time $L$.

\begin{proposition}[Cone identity and an average of squares]
  \label{P:point-source-cone}
  Let $u_{\epsilon,\tau}^{\rm full}$ be the unmasked source solution with
  coefficient $gY_{\rm fast}$ and data $(1,0)$. At the fixed parameters, almost
  surely,
  \[
    v_{\epsilon,\tau}=u_{\epsilon,\tau}^{\rm full}
    \quad\hbox{on }\mathcal{C}=
    \{(t,x):0\le t\le L,\ |x|\le1+L-t\}.
  \]
  Consequently,
  \begin{equation}\label{E:point-source-square-average}
    \E\|\operatorname{pr}_1P_E^{\epsilon,\tau}(L,0)z\|_{L^2(D)}^2
    =v_d\,\E|u_{\epsilon,\tau}^{\rm full}(L,0)|^2<\infty .
  \end{equation}
\end{proposition}

\begin{proof}
  We compare the compact-data and constant-data source equations on one
  causally closed cone. Their free terms agree there by total finite
  propagation, and iteration of the homogeneous difference equation gives
  equality before taking the terminal spatial square average.


  \medskip
  \noindent\textit{Step 1.~}We identify the common source law and the causally
  closed cone.
  Proposition~\ref{P:masked-source-moments}, applied through the exact source
  law of Lemma~\ref{L:point-fast-sources}, gives the locally continuous
  measurable full solution and its $L^2$ point Picard sum. The pathwise cone
  argument uses the realized pullback field; the point second moments depend
  only on its Gaussian covariance through those Bochner chain sums. We place no
  global unmasked field in $E$ or $F$. The cone is causally closed: a
  predecessor at distance at most $t-s$ has norm at most $1+L-s$. Its spatial
  projection lies a positive distance inside the plateau of both $\psi$ and $m$.


  \medskip
  \noindent\textit{Step 2.~}We compare both initial phase components on the
  cone.
  The total initial difference $(\psi-1,0)$ is represented in $E$ plus constants
  by $(\psi,0)+(-1,0)$. Every entry of the full free phase kernel is a compactly
  supported distribution in $B_t$, and the top-left is $\partial_tG_d(t)$. Since
  $|x|+t \le 1+L = R-1$ on the cone, its entire translated kernel support lies
  in the region where this total datum vanishes.
  Lemma~\ref{L:constant-phase-free} therefore gives
  $\cos(t\sqrt{-\Delta})\psi(x) = 1$ there. Testing the $E$ and constant
  summands separately would not give this total finite-propagation statement.


  \medskip
  \noindent\textit{Step 3.~}We prove uniqueness of the two restrictions on the
  cone.
  The two mild equations have the same coefficient on the cone. For their
  difference $w$, causal closure gives the homogeneous Green integral equation
  there. Put $H = \sup_{\mathcal{C}}|w|$ and $M = \sup_{\mathcal{C}}|gY_{\rm
  fast}|$, both finite sample constants. Iteration with the positive Green
  measures gives $|w(t,x)| \le HM^nt^{2n}/(2n)!$. Indeed, the mass recursion is
  $f_0 = 1$, $f_{n+1}(t) = \int_0^t(t-r)f_n(r)\ud r = t^{2n+2}/(2n+2)!$. The
  factorial ratio forces this bound to zero on the fixed cone. No positivity of
  the sample coefficient or expectation of a global source supremum is required.


  \medskip
  \noindent\textit{Step 4.~}We take the terminal average of squares.
  At time $L$ the cone contains $D$. Tonelli and spatial stationarity of the
  unmasked single-source law give~\eqref{E:point-source-square-average}.
  Finiteness is the full source point-moment bound. This is an integral of
  squares, not a squared average or a covariance comparison with a masked model.
  Events are for each fixed $\delta,\epsilon,\tau$. This proves
  Proposition~\ref{P:point-source-cone}.
\end{proof}

Agreement on the terminal observation set alone would not suffice:
both initial phase components must agree on its backward cone.
No-go~\ref{P:nogo-31} gives a smooth free-wave counterexample when
the velocity guard is omitted, and a separate phase example explaining
why operator growth alone does not imply growth for every initial state.


Proposition~\ref{P:point-source-cone} holds at fixed positive
$\epsilon,\tau$ and for the scaled covariance $\Gamma_\delta$. We now remove the
regularization and undo the scaling to reach the original point moment
$\E|u_\Gamma(T_\delta,0)|^2$.

\begin{proposition}[Exact original point transfer and polynomial cost]
  \label{P:original-point-transfer}
  For every original model, including zero covariance,
  \begin{align}
    \E|u_\Gamma(T_\delta,0)|^2
    &=v_d^{-1}\E\|\operatorname{pr}_1P_\delta z\|_{L^2(D)}^2,
    \label{E:original-point-transfer}\\
    &\le v_d^{-1}\|\psi\|_{H^{1/2}}^2\E\|P_\delta\|^2,\notag\\
    \limsup_{T\to\infty}T^{-\kappa}\log\E|u_\Gamma(T,x)|^2
    &\le\limsup_{\delta\downarrow0}\delta\log\E\|P_\delta\|^2
    \quad\hbox{for every fixed }x .
    \label{E:original-point-rate-reduction}
  \end{align}
\end{proposition}

\begin{proof}
  We remove source regularization at fixed dilation, then use the canonical
  law scaling to identify the original point moment. The compact initial datum
  costs only a polynomial factor, which disappears after logarithmic
  normalization.


  \medskip
  \noindent\textit{Step 1.~}We pass the source square identity to the original
  operator endpoint.
  Assume nonzero covariance until the last paragraph, and fix $\delta$.
  Equation~\eqref{E:point-source-flow-limit} gives convergence of the source
  displacement in $L^2(\Omega;L^2(D))$, since $\|\operatorname{pr}_1e\|_{L^2(D)}
  \le \|e\|_E$. Expected squared norms converge by Cauchy--Schwarz applied to
  their difference times their sum. Separately, the full source point-moment
  limit \eqref{E:masked-unmasked-moment-limit} at covariance $\Gamma_\delta$
  gives convergence of the right side of~\eqref{E:point-source-square-average}
  to $v_d\E|u_{\Gamma_\delta}(L,0)|^2$. Both cutoffs $(\epsilon,\tau/\delta)$
  tend to zero at fixed $\delta$. Thus, the two limits in the equal prelimit
  moments give
  \[
    \E|u_{\Gamma_\delta}(L,0)|^2
    =v_d^{-1}\E\|\operatorname{pr}_1P_\delta z\|_{L^2(D)}^2 .
  \]


  \medskip
  \noindent\textit{Step 2.~}We identify the original moment by the canonical
  time change.
  The law normalization is the consequence of
  Lemma~\ref{L:exact-canonical-scaling} and the canonical-series time change
  proved in Proposition~\ref{P:raw-block-point-comparison}. Explicitly, for
  $\delta = T^{-\kappa}$ and $r_T = T\delta$, $V_\delta$ has the law of
  $u_\Gamma(Tt,r_Tx)$ and equation \eqref{E:exact-fast-equation}. At positive
  source widths, $U(t,x) = V_\delta(\delta t,x)$ has coefficient
  $\sqrt\delta\,Y_{\epsilon,\tau}(\delta t,x) = gY_{\rm fast}(t,x)$. The
  fixed-order and absolutely convergent canonical-series argument passes this
  time change to the point moment at fixed $\delta$. Hence,
  $\E|u_\Gamma(T_\delta,0)|^2 = \E|u_{\Gamma_\delta}(L,0)|^2$. This is a
  law/moment identity, not a pathwise identity on the same noise. It preserves
  all mixed, unequal, self and empty Stratonovich chains and the
  Green-before-noise convention in $d = 3$. The shrinking $r_T$ is not the mask
  radius $R$; no spatial Jacobian occurs at point zero. The phase/operator
  inequality now proves~\eqref{E:original-point-transfer}.


  \medskip
  \noindent\textit{Step 3.~}We control the initial-data norm and its
  logarithmic cost.
  Parseval and the inhomogeneous Bessel weights give, for $R \ge 1$,
  \[
    \|\psi\|_{H^{1/2}}^2\le\|\psi\|_{H^1}^2
    =R^d\|\eta\|_2^2+R^{d-2}\|\nabla\eta\|_2^2
    \le C_\eta R^d,\quad C_\eta=\|\eta\|_2^2+\|\nabla\eta\|_2^2>0.
  \]
  This includes the negative $d-2$ power in $d = 1$. For
  $\operatorname{DataPref}_\delta = v_d^{-1}C_\eta(\delta^{-1}+2)^d$,
  \[
    \delta\log\operatorname{DataPref}_\delta
    =\delta\log(C_\eta/v_d)+d\delta[\log(1+2\delta)-\log\delta]\to0 .
  \]
  Since $R\ge1$, the factor $R^{d-2}$ is at most $R^d$ even when $d-2$ is
  negative. The squared initial norm is therefore bounded by the same volume
  power in all three dimensions. Substituting $R=\delta^{-1}+2$ gives only a
  logarithmic factor of order $\log(1/\delta)$, whose product with $\delta$
  tends to zero.
  The nontrivial term follows by $t = -\log\delta$ and $te^{-t} \to 0$, since
  $e^t \ge t^2/2$. The original point moment is at least one: its nonnegative
  full pairing expansion contains the unit empty pair of chains. It is finite
  by~\eqref{E:original-point-transfer} and~\eqref{E:point-fixed-finiteness}.
  Thus, logs are defined without subtracting infinite quantities. Since $\alpha
  < 2 < 3$, $a < 1$, we have $\kappa > 1$ and $\delta\mapsto T_\delta$ is a
  decreasing bijection $(0,1) \to (1,\infty)$ with inverse $\delta =
  T^{-\kappa}$. Taking logs and this change of parameter proves
  \eqref{E:original-point-rate-reduction} at zero. Original canonical spatial
  stationarity, proved in Proposition~\ref{P:raw-block-original-lower}, extends
  it to every fixed $x$; no stationarity of the masked operator is assumed.


  \medskip
  \noindent\textit{Step 4.~}We verify the free zero-covariance case directly.
  For zero covariance, $u = 1$, all operator paths vanish and $P_\delta =
  S_E(L)$. The same total free cone argument gives the spatial-square identity.
  Lemma~\ref{L:symplectic-free} gives $\|S_E(L)\| \ge 1$, and its Fourier matrix
  gives $\|S_E(L)\| \le C(1+L)$. Thus, $0 \le \delta\log\|S_E(\delta^{-1})\|^2
  \to 0$ and both claims hold directly without a nonzero-noise import. This
  proves Proposition~\ref{P:original-point-transfer}.
\end{proof}

\begin{figure}
  \centering
  \begin{tikzpicture}[x=1cm,y=1cm]
    \fill[hamblue!4] (-5.4,0) rectangle (5.4,3.4);
    \fill[white] (-2.7,0) rectangle (2.7,3.4);
    \fill[hamteal!10] (-2.5,0)--(0,2.5)--(2.5,0)--cycle;
    \draw[ham/arrow] (-5.7,0)--(5.7,0) node[right] {$y$};
    \draw[ham/arrow] (0,-.25)--(0,3.55) node[above] {$s$};
    \draw[thick,hamteal] (-2.5,0)--(0,2.5)--(2.5,0);
    \draw[dashed,hamblue] (-2.7,0)--(-2.7,3.25);
    \draw[dashed,hamblue] (2.7,0)--(2.7,3.25);
    \draw[hamblue] (-5.4,0)--(-5.4,3.25);
    \draw[hamblue] (5.4,0)--(5.4,3.25);
    \fill (0,2.5) circle (1.7pt);
    \node[anchor=west,fill=white,inner sep=2pt] at (.15,2.55) {$(T,x)$};
    \node[ham/note,text width=4.2cm] at (0,.65)
    {$|y-x|\le T-s$\\data: $(u,\partial_tu)$};
    \node[font=\small,anchor=north] at (-2.7,-.1) {$-R$};
    \node[font=\small,anchor=north] at (2.7,-.1) {$R$};
    \node[ham/note,fill=white] at (0,3.05) {mask equals one on $B_R$};
    \node[font=\small,anchor=north] at (-5.4,-.1) {$-2R$};
    \node[font=\small,anchor=north] at (5.4,-.1) {$2R$};
  \end{tikzpicture}
  \caption{A spatial cross-section of finite propagation. The entire backward
    cone is chosen inside the region where the mask equals one; the mask
    vanishes outside $\overline{B}_{2R}$. The drawing is schematic and does not
    prescribe a value of $R/T$. Shading indicates the domain of dependence, not
    a Green density; in dimension three the Green measure is supported on a
    sphere. In the point-moment comparison, $L=\delta^{-1}$, $R=L+2$ and $x\in
    B_1$. Both displacement and velocity must be passed between time intervals.
    See \ProofRef{Lemma}{L:raw-block-restart} and
    \ProofRef{Proposition}{P:original-point-transfer}.}
  \label{Fig:light-cone}
\end{figure}

